\chapter{Risco de mercado: VaR, CVaR e drawdown}
\label{cap:risco}

\abertura{Uma carteira de ações com volatilidade de 25\% ao ano é arriscada, mas esse número não diz ao gestor quantos reais ele pode perder amanhã, nem de que tamanho é o prejuízo nos dias realmente ruins, nem quanto tempo a carteira pode passar abaixo do seu melhor momento. Este capítulo apresenta as três medidas que respondem a essas perguntas, o VaR, o \ingles{expected shortfall} e o \ingles{drawdown}, e mostra como verificar se os números que elas produzem merecem confiança.}

\begin{objetivos}
  \item definir o VaR como um quantil da distribuição de perdas e interpretá-lo sem exagerar o que ele diz;
  \item calcular o VaR de uma carteira pelos métodos paramétrico, histórico e de Monte Carlo, e explicar por que os resultados diferem;
  \item testar um modelo de VaR contando violações e aplicando o teste de Kupiec;
  \item calcular o \ingles{expected shortfall} e explicar por que ele substituiu o VaR na regulação bancária;
  \item construir a série de \ingles{drawdown} de uma carteira e extrair dela o drawdown máximo, a duração e o tempo de recuperação;
  \item reconhecer o tamanho do erro de estimação em qualquer medida de cauda.
\end{objetivos}

\section{Risco como perda em reais}

A volatilidade, que você aprendeu a calcular no Capítulo~\ref{cap:retornos} e a modelar no Capítulo~\ref{cap:series}, mede a dispersão dos retornos em torno da média. Ela tem duas limitações como medida de risco. Trata da mesma forma um desvio para cima e um desvio para baixo, embora ninguém considere um ganho inesperado um problema. E resume a distribuição inteira em um número que só descreve bem as caudas quando os retornos são normais, o que o Capítulo~\ref{cap:estatistica} mostrou não ser o caso: duas carteiras com a mesma volatilidade podem ter piores dias de tamanhos muito diferentes.

As medidas deste capítulo olham diretamente para o lado das perdas e são expressas na unidade que interessa a quem responde pelo dinheiro. Para isso é preciso fixar uma convenção de sinal, que será mantida até o fim do capítulo.

\begin{definicao}{Perda}
Para uma carteira de valor $V$ hoje e retorno simples $R_p$ no horizonte considerado, a \termo{perda} é
\begin{equation}
  L = -V R_p .
  \label{eq:c05-perda}
\end{equation}
Perdas são números positivos e ganhos são perdas negativas.
\end{definicao}

Todos os exemplos com dados usam a mesma carteira: R\$~1 milhão dividido igualmente entre PETR4, VALE3, ITUB4, WEGE3 e ABEV3. O Código~\ref{cod:c05-carteira} monta a série de perdas diárias que essa carteira teria sofrido em cada dia do histórico, e os demais códigos do capítulo a importam.

\lstinputlisting[style=py, caption={A carteira do capítulo e sua série de perdas (\arq{cap05/carteira.py})}, label=cod:c05-carteira, lastline=22]{scripts/cap05/carteira.py}
\begin{lstlisting}[style=saida]
2471 dias, de 02/01/2015 a 30/12/2024
perda média: R$ -905
desvio-padrão: R$ 16,026
maior perda: R$ 150,307 em 12/03/2020
retornos de mais de um pregão: 30 (3 de 2015 a 2024)
\end{lstlisting}

A saída, impressa pelo bloco final do arquivo, resume os dez anos de 2015 a 2024, que servem de amostra de estimação no restante do capítulo. A perda média é negativa porque a carteira ganhou, em média, R\$~905 por dia. A série de perdas diz quanto a carteira de \emph{hoje}, com os pesos de hoje, perderia se o mercado repetisse cada dia do passado, que é o que interessa para medir o risco de amanhã; ela é diferente do resultado de quem aplicou R\$~1 milhão em 2015 e deixou render.

A base de preços tem valores ausentes nos pregões em que um papel não teve negócio registrado. Essas datas saem, e o retorno seguinte cobre mais de um pregão: o de 17 de agosto de 2011 reúne 14 pregões, porque PETR4 ficou sem preço na base. Há 30 retornos assim de 2010 a 2025, quase todos em 2011 e 2012, e três entre 2015 e 2024. Eles entram na \cod{cota}, usada no drawdown da Seção~\ref{sec:c05-drawdown}, que precisa de todo o caminho, e ficam fora da série de perdas diárias, do VaR e do backtest, em que apareceriam como perdas de um dia que não ocorreram.

\section{Value at Risk: um quantil das perdas}

A Figura~\ref{fig:c05-perdas} mostra o histograma das 2471 perdas diárias. Em 124 desses dias, 5\% do total, a carteira perdeu mais de R\$~22,9 mil. Esse valor de corte é o \ingles{Value at Risk} de 95\% da carteira para um dia, estimado da forma mais simples possível. Em 25 dias, 1\% do total, a perda passou de R\$~37,6 mil: é o VaR de 99\%.

\figura[t]{cap05/perdas}{Perdas diárias da carteira de R\$~1 milhão entre 2015 e 2024. O VaR de 95\% é o valor que separa os 5\% piores dias (em ocre) dos demais. As barras das duas bordas reúnem os dias além de R\$~80 mil, de ganho ou de perda. O histograma mostra onde a cauda começa; o VaR nada diz sobre até onde ela vai.}{fig:c05-perdas}

\begin{definicao}{Value at Risk}
\index{VaR|see{Value at Risk}}O \termo{Value at Risk} de nível de confiança $\alpha$ é o quantil $\alpha$ da distribuição da perda $L$ em um horizonte fixado:
\begin{equation}
  \VaR_\alpha(L) = \inf\{\ell \in \R : P(L \le \ell) \ge \alpha\} .
  \label{eq:c05-var}
\end{equation}
Se $L$ tem distribuição contínua, $\VaR_\alpha$ é o valor que a perda supera com probabilidade exatamente $1-\alpha$.
\end{definicao}

A definição segue \citet{mcneil2015}. Em termos de retornos, a Equação~\eqref{eq:c05-perda} dá $\VaR_\alpha(L) = -V q_{1-\alpha}(R_p)$, em que $q_{1-\alpha}$ é o quantil $1-\alpha$ dos retornos: o VaR de 95\% é o quinto percentil dos retornos, com o sinal trocado e multiplicado pelo valor da carteira.\nota{Muitos textos e os scripts da aula informam o VaR como um retorno negativo, por exemplo $-2,3\%$. É a mesma medida com outra convenção de sinal. Ao ler um relatório, descubra primeiro qual convenção ele usa.}

Um VaR só tem significado acompanhado de duas escolhas. O \termo{nível de confiança} $\alpha$ costuma ser 95\% ou 99\%; quanto maior, mais para dentro da cauda se olha e, como se verá na Seção~\ref{sec:c05-cauda}, menos dados sustentam a estimativa. O \termo{horizonte} é o intervalo em que a perda é medida, normalmente um dia para mesas de operação e dez dias úteis na regulação bancária tradicional. A frase completa é sempre do tipo ``o VaR de um dia a 99\% é de R\$~37,6 mil''.

\begin{armadilha}{O VaR não é a perda máxima}
O VaR de 99\% não diz que a carteira perde no máximo R\$~37,6 mil. Diz que em um dia a cada cem, em média, ela perde \emph{mais} do que isso, sem informar quanto mais. Na amostra da Figura~\ref{fig:c05-perdas} houve três dias com perda acima de R\$~130 mil, mais de três vezes o VaR de 99\%.
\end{armadilha}

O que muda de um método de cálculo para outro é a forma de obter a distribuição de $L$: supor uma família de distribuições, usar o histograma do passado ou simular.

\section{VaR paramétrico}
\label{sec:c05-param}

\subsection{Sob normalidade}

Se a perda tem distribuição normal de média $\mu_L$ e desvio-padrão $\sigma_L$, seus quantis são conhecidos e o VaR sai de uma conta.

\begin{resultado}{VaR normal}
Se $L \sim N(\mu_L, \sigma_L^2)$, então
\begin{equation}
  \VaR_\alpha(L) = \mu_L + \sigma_L\, z_\alpha ,
  \label{eq:c05-var-normal}
\end{equation}
em que $z_\alpha$ é o quantil $\alpha$ da normal padrão: $z_{0,95} = 1,645$ e $z_{0,99} = 2,326$.
\end{resultado}

Com $\mu_L = -V\mu_p$ e $\sigma_L = V\sigma_p$, a fórmula em retornos é $\VaR_\alpha = V(z_\alpha \sigma_p - \mu_p)$. Em horizonte de um dia a média é pequena diante do desvio-padrão (na nossa carteira, R\$~905 contra R\$~16 mil) e muito mal estimada, e por isso é prática comum fixá-la em zero.

Para uma carteira de vários ativos com vetor de pesos $\vetor{w}$, vetor de médias $\symbf{\mu}$ e matriz de covariância $\symbf{\Sigma}$ dos retornos, a Equação~\eqref{eq:ret-carteira} do Capítulo~\ref{cap:retornos} implica
\begin{equation}
  \mu_p = \vetor{w}\T \symbf{\mu} , \qquad \sigma_p = \sqrt{\vetor{w}\T \symbf{\Sigma}\, \vetor{w}} .
  \label{eq:c05-sigma-p}
\end{equation}
A segunda expressão carrega todo o efeito da diversificação, e o Capítulo~\ref{cap:portfolio} a estuda em detalhe. Aqui basta um exemplo para ver o que ela faz com o VaR.

\begin{exemplo}{VaR de dois ativos, à mão}
Uma carteira de R\$~100 mil tem 60\% no ativo A, com volatilidade diária de 2\%, e 40\% no ativo B, com volatilidade diária de 3\%. A correlação entre eles é $0,3$ e as médias são tomadas como zero. A variância da carteira é
\[
  \sigma_p^2 = 0,6^2 \cdot 0,02^2 + 0,4^2 \cdot 0,03^2 + 2 \cdot 0,6 \cdot 0,4 \cdot 0,3 \cdot 0,02 \cdot 0,03 = 0,0003744 ,
\]
de modo que $\sigma_p = 1,935\%$. O VaR de um dia a 95\% é $1,645 \times 0,01935 \times 100\,000 = \text{R\$~}3\,183$.

Isoladas, a posição em A tem VaR de $1,645 \times 0,02 \times 60\,000 = \text{R\$~}1\,974$ e a posição em B, por coincidência, o mesmo valor: $1,645 \times 0,03 \times 40\,000 = \text{R\$~}1\,974$. A soma é R\$~3\,948. Os R\$~765 de diferença são o ganho de diversificação, que só desapareceria se a correlação fosse igual a 1.
\end{exemplo}

O Código~\ref{cod:c05-param} faz a mesma conta para a carteira de cinco ações, agora com a média estimada.

\needspace{14\baselineskip}
\lstinputlisting[style=py, caption={VaR paramétrico normal a partir de $\vetor{w}$ e $\symbf{\Sigma}$ (\arq{cap05/l01\_parametrico.py})}, label=cod:c05-param]{scripts/cap05/l01_parametrico.py}
\begin{lstlisting}[style=saida]
mu_L = -905   sigma_L = 16,026
VaR normal 95%: z = 1.645  R$ 25,456
VaR normal 99%: z = 2.326  R$ 36,377
soma dos VaR 99% isolados (média zero): R$ 53,192
VaR 99% da carteira (média zero):       R$ 37,282
\end{lstlisting}

O desvio-padrão de R\$~16.026 obtido por $\vetor{w}\T\symbf{\Sigma}\vetor{w}$ é idêntico ao da série de perdas do Código~\ref{cod:c05-carteira}, como tem de ser: são duas maneiras de calcular a mesma variância amostral. As duas últimas linhas mostram que a diversificação entre as cinco ações reduz o VaR em 30\% em relação à soma das posições isoladas.

\subsection{Com a t de Student}

O VaR normal de 95\% (R\$~25.456) ficou acima do valor de R\$~22,9 mil lido no histograma, e o de 99\% (R\$~36.377) ficou abaixo dos R\$~37,6 mil. O desencontro tem a assinatura das caudas pesadas descritas no Capítulo~\ref{cap:estatistica}: em relação à normal de mesma variância, os retornos reais têm mais massa perto do centro e nos extremos, e menos nos ``ombros''. Uma alternativa paramétrica que acomoda isso é a t de Student.

Uma perda com distribuição t de Student de $\nu$ graus de liberdade, locação $m$ e escala $s$ tem
\begin{equation}
  \VaR_\alpha(L) = m + s\, t_{\nu,\alpha} ,
  \label{eq:c05-var-t}
\end{equation}
em que $t_{\nu,\alpha}$ é o quantil $\alpha$ da t padrão. Há um cuidado de escala: $s$ não é o desvio-padrão. Uma variável com distribuição t padrão e $\nu > 2$ graus de liberdade tem variância $\nu/(\nu-2)$, maior que 1, de modo que o desvio-padrão da perda é $s\sqrt{\nu/(\nu-2)}$. Quem parte da média e do desvio-padrão amostrais precisa usar $m = \mu_L$ e $s = \sigma_L\sqrt{(\nu-2)/\nu}$.

\needspace{7\baselineskip}
\lstinputlisting[style=py, caption={VaR com t de Student ajustada por máxima verossimilhança (\arq{cap05/l02\_student.py})}, label=cod:c05-student]{scripts/cap05/l02_student.py}
\begin{lstlisting}[style=saida]
nu = 4.48   m = -957   s = 11,660
desvio-padrão implícito: 15,672   amostral: 16,026
VaR t 95%: quantil = 2.068  ajuste: R$ 23,156  momentos: R$ 23,752
VaR t 99%: quantil = 3.535  ajuste: R$ 40,258  momentos: R$ 41,240
VaR t 99% sem corrigir a escala: R$ 55,688
\end{lstlisting}

A função \cod{stats.t.fit} estima de uma vez $\nu$, $m$ e $s$ por máxima verossimilhança. O ajuste encontra $\nu = 4,48$, caudas bem mais pesadas que as da normal, e uma escala de R\$~11.660, que corresponde a um desvio-padrão de R\$~15.672, um pouco menor que o amostral. Por isso as duas colunas da saída diferem: ``ajuste'' aplica a Equação~\eqref{eq:c05-var-t} com o $m$ e o $s$ estimados, e ``momentos'' usa o mesmo $\nu$ com $m = \mu_L$ e $s = \sigma_L\sqrt{(\nu-2)/\nu}$, o que dá um VaR cerca de 2,5\% maior (R\$~41.240 contra R\$~40.258 a 99\%). As duas são leituras legítimas do mesmo modelo, e a diferença entre elas é erro de estimação. A última linha mostra um erro de outra natureza, o de usar o desvio-padrão amostral diretamente como escala: o VaR de 99\% sobe para R\$~56 mil, um exagero de quase 40\%.

\figura[t]{cap05/normal_t}{Quantis da perda, em desvios-padrão, em função do nível de confiança. A normal e a t de mesma variância se cruzam perto de 97\%: abaixo desse nível a normal dá o VaR maior, acima dele a t dá o maior. Os pontos são os quantis empíricos da carteira; os dois últimos se apoiam em cinco e em duas ou três observações.}{fig:c05-normal-t}

A Figura~\ref{fig:c05-normal-t} explica os números. As duas curvas têm a mesma variância e se cruzam perto do nível de 97\%. A 95\% a hipótese normal é a mais conservadora das duas; a 99\% ela subestima a perda, e a diferença cresce rapidamente depois disso. Os quantis empíricos da carteira ficam junto da curva da t, e abaixo da normal, até 97,5\%, entre as duas a 99\% e a 99,5\%, e acima das duas a partir de 99,8\%.

\subsection{Mudança de horizonte e a raiz do tempo}

Se os retornos diários forem independentes, com a mesma distribuição normal e média zero, a perda acumulada em $h$ dias tem desvio-padrão $\sigma_L\sqrt{h}$ (é a regra da raiz do tempo da Seção~\ref{sec:anualizar}), e
\begin{equation}
  \VaR_\alpha^{(h\text{ dias})} = \sqrt{h}\;\VaR_\alpha^{(1\text{ dia})} .
  \label{eq:c05-raiz}
\end{equation}
Para a carteira do capítulo, o VaR normal de dez dias a 99\% com média zero seria $\sqrt{10} \times 37.282$, cerca de R\$~117,9 mil.

A regra é uma convenção de escala, e as hipóteses por trás dela falham. O script \arq{cap05/horizonte.py} compara $\sqrt{h}$ vezes o VaR histórico de um dia com o quantil das perdas acumuladas em janelas sobrepostas de $h$ dias. A 99\%, a regra dá R\$~84,0 mil contra R\$~74,6 mil observados em cinco dias, R\$~118,8 mil contra R\$~121,0 mil em dez dias e R\$~172,1 mil contra R\$~198,7 mil em 21 dias; a 95\% e dez dias, R\$~72,5 mil contra R\$~60,8 mil. Duas forças atuam em sentidos opostos. A média positiva e a autocorrelação levemente negativa dos retornos diários desta carteira ($-0,08$) contribuem para que a perda acumulada cresça menos que $\sqrt{h}$ no miolo da distribuição; o agrupamento de volatilidade (Capítulo~\ref{cap:series}) atua no sentido de engordar a cauda extrema. Leia os números de 99\% com reserva: as janelas se sobrepõem, e 20 das 25 janelas de 21 dias acima do quantil terminam em março ou abril de 2020. A comparação nesse nível se apoia em um único episódio.

\section{VaR histórico}

O método histórico, ou de simulação histórica, não supõe família alguma de distribuições. Ele toma as perdas que a carteira atual teria sofrido em cada dia de uma janela do passado e usa o quantil empírico:
\begin{equation}
  \widehat{\VaR}_\alpha = \text{quantil empírico de nível } \alpha \text{ de } \{L_1, \dots, L_n\} .
\end{equation}
Foi o que se fez na Figura~\ref{fig:c05-perdas}.

\needspace{18\baselineskip}
\lstinputlisting[style=py, caption={VaR histórico e sua dependência da janela (\arq{cap05/l03\_historico.py})}, label=cod:c05-hist]{scripts/cap05/l03_historico.py}
\begin{lstlisting}[style=saida]
VaR histórico 95%: R$ 22,914
VaR histórico 99%: R$ 37,561
VaR histórico 99%, 2015 a 2019: R$ 37,390
VaR histórico 99%, 2020 a 2024: R$ 37,603
VaR histórico 99%, 2024 a 2024: R$ 20,168
\end{lstlisting}

Com 2471 observações, o quantil de 99\% cai entre a 25ª e a 26ª maior perda, e \cod{np.quantile} interpola entre as duas. A vantagem do método é evidente. Assimetria, caudas pesadas e a dependência entre os ativos nos dias ruins entram no cálculo sem que seja preciso modelá-las, porque já estão nos dados.

O preço é a dependência da janela, visível nas três últimas linhas da saída. Com apenas os 249 dias de 2024, um ano calmo, o VaR de 99\% cai para R\$~20 mil, pouco mais da metade do valor obtido com dez anos. O método só conhece o que aconteceu dentro da janela e supõe, implicitamente, que o futuro próximo será sorteado da mesma distribuição. Uma janela curta reage depressa a mudanças de regime, mas apoia o quantil em duas ou três observações; uma janela longa dá mais observações, mas mistura regimes de volatilidade muito diferentes e reage devagar.

\figura[t]{cap05/fantasma}{O efeito fantasma. Com janela de 250 dias, o VaR histórico de 99\% triplica em março de 2020 e permanece no patamar novo por um ano, até despencar no dia em que as perdas daquele mês saem da janela. Com 1000 dias, o mesmo choque mal altera o quantil.}{fig:c05-fantasma}

A Figura~\ref{fig:c05-fantasma} mostra a consequência mais estranha dessa mecânica, conhecida como \termo{efeito fantasma}. Com 250 observações, o quantil de 99\% fica entre a terceira e a quarta maior perda da janela. Em março de 2020 três perdas acima de R\$~130 mil entram na janela e o VaR salta de R\$~32 mil para R\$~105 mil. Ele fica nesse patamar por um ano, embora o mercado tenha se acalmado em poucos meses, e em 16 de março de 2021 cai de R\$~105 mil para R\$~77 mil de um dia para o outro, sem que nada tenha acontecido no mercado: a perda de 9 de março de 2020 deixou de estar entre os 250 pregões mais recentes e saiu da conta. A janela móvel dá o mesmo peso a todas as observações que contém e peso zero às demais, e essa descontinuidade produz o degrau.

\section{VaR por simulação de Monte Carlo}

\index{Monte Carlo}O terceiro caminho consiste em escrever um modelo para os retornos dos ativos, sortear dele um grande número de cenários, calcular a perda da carteira em cada cenário e tomar o quantil empírico das perdas simuladas. É o método histórico aplicado a um passado artificial, do tamanho que se quiser.

Para sortear retornos de vários ativos é preciso reproduzir a covariância entre eles. O instrumento é a \termo{decomposição de Cholesky}: toda matriz de covariância positiva definida pode ser escrita como $\symbf{\Sigma} = \vetor{C}\vetor{C}\T$, com $\vetor{C}$ triangular inferior. Se $\vetor{z}$ é um vetor de normais padrão independentes, então $\vetor{x} = \vetor{C}\vetor{z}$ tem covariância
\begin{equation}
  \Cov(\vetor{C}\vetor{z}) = \vetor{C}\,\Cov(\vetor{z})\,\vetor{C}\T = \vetor{C}\vetor{C}\T = \symbf{\Sigma} ,
\end{equation}
pois $\Cov(\vetor{z})$ é a identidade. O Código~\ref{cod:c05-mc} usa essa ideia para gerar 200 mil cenários de um dia.

\lstinputlisting[style=py, caption={VaR por Monte Carlo com choques correlacionados (\arq{cap05/l04\_montecarlo.py})}, label=cod:c05-mc]{scripts/cap05/l04_montecarlo.py}
\begin{lstlisting}[style=saida]
Monte Carlo normal  VaR 95%: R$ 25,381   VaR 99%: R$ 36,319
Monte Carlo t       VaR 95%: R$ 23,730   VaR 99%: R$ 41,533
\end{lstlisting}

Cada linha de \cod{Z} é um cenário e \cod{Z @ C.T} aplica $\vetor{C}$ a todos de uma vez. A segunda metade do código troca a normal por uma t multivariada: dividir um vetor normal pela raiz de uma qui-quadrado independente, dividida por seus graus de liberdade, produz uma t, e o fator $\sqrt{(\nu-2)/\nu}$ é a mesma correção de escala da Equação~\eqref{eq:c05-var-t}, agora para que a covariância continue sendo $\symbf{\Sigma}$. O valor de $\nu$ é o estimado no Código~\ref{cod:c05-student}.

\begin{table}[b]
\caption{VaR de um dia da carteira de R\$~1 milhão por seis caminhos, com os mesmos dados de 2015 a 2024 (valores em reais, saídas dos Códigos~\ref{cod:c05-param} a~\ref{cod:c05-mc})}
\label{tab:c05-comparacao}
\begin{tabularx}{\textwidth}{>{\raggedright\arraybackslash}X r r}
\toprule
\cab{Método} & \cab{VaR 95\%} & \cab{VaR 99\%} \\
\midrule
Paramétrico normal & 25.456 & 36.377 \\
Paramétrico t, máxima verossimilhança & 23.156 & 40.258 \\
Paramétrico t, momentos amostrais & 23.752 & 41.240 \\
Histórico & 22.914 & 37.561 \\
Monte Carlo normal, 200 mil cenários & 25.381 & 36.319 \\
Monte Carlo t, 200 mil cenários & 23.730 & 41.533 \\
\bottomrule
\end{tabularx}
\end{table}

A Tabela~\ref{tab:c05-comparacao} reúne os resultados. O Monte Carlo normal reproduz o paramétrico normal a menos de 0,3\%, e o Monte Carlo t reproduz a menos de 0,7\% o paramétrico t pelos momentos, com o qual compartilha a média e a covariância amostrais. Isso não é coincidência. Para uma carteira que é combinação linear dos ativos, sob um modelo normal ou t multivariado, a distribuição da perda é conhecida em forma fechada, e a simulação apenas a aproxima com ruído. A diferença que importa na tabela é entre hipóteses de distribuição, e não entre métodos de cálculo: a 99\%, os modelos com t de Student pedem de 11\% a 14\% a mais que os normais.

A simulação acrescenta, além disso, um erro próprio. Repetindo 200 vezes o Monte Carlo normal (script \arq{cap05/mc\_erro.py}), o VaR de 99\% estimado com mil cenários varia de R\$~32 mil a R\$~41 mil entre rodadas, com desvio-padrão de 5\% do valor exato; com 100 mil cenários o desvio-padrão cai para 0,5\%. Como o erro de simulação cai com $1/\sqrt{N}$, cada casa decimal a mais exige cem vezes mais cenários.

O Monte Carlo compensa esse custo quando não existe fórmula fechada. O caso clássico é o de carteiras com opções, cujo valor não varia linearmente com o preço do ativo: a perda da carteira deixa de ser normal mesmo que os retornos dos ativos sejam, e é preciso reavaliar cada posição em cada cenário \citep{hull2018, jorion2007}. Os outros casos são modelos dinâmicos, como simular dez dias de um GARCH (Capítulo~\ref{cap:series}) para obter o VaR de dez dias sem recorrer à raiz do tempo, e regras que dependem da trajetória, como uma ordem de \ingles{stop}, que encerra a posição quando a perda atinge um limite. Em todos eles vale a mesma ressalva: a simulação reproduz fielmente o modelo que recebeu, inclusive os defeitos dele.

\section{Backtest do VaR: contar violações}
\label{sec:c05-backtest}

Um modelo de VaR faz uma afirmação que pode ser conferida. Se o VaR de 99\% estiver certo, a perda do dia seguinte deve superá-lo em 1\% dos dias. Cada dia em que isso ocorre é uma \termo{violação} (ou exceção). O \ingles{backtest} de um modelo de VaR consiste em recalcular o VaR a cada dia usando só os dados disponíveis até a véspera, comparar com a perda realizada e contar.

Violações de menos também são defeito: indicam um VaR alto demais, que imobiliza capital sem necessidade. O teste formal mais usado para a contagem é o de \citet{kupiec1995}.\index{teste de Kupiec}

\begin{resultado}{Teste de proporção de falhas de Kupiec}
Sejam $T$ o número de dias do backtest, $x$ o número de violações, $p = 1-\alpha$ a probabilidade de violação prometida pelo modelo e $\hat{p} = x/T$ a frequência observada. Sob a hipótese nula de que as violações são independentes e têm probabilidade $p$, a estatística
\begin{equation}
  LR = -2 \ln\!\left[\frac{(1-p)^{T-x}\, p^{x}}{(1-\hat{p})^{T-x}\, \hat{p}^{x}}\right]
  \label{eq:c05-kupiec}
\end{equation}
tem, para $T$ grande, distribuição aproximadamente qui-quadrado com um grau de liberdade. Rejeita-se o modelo ao nível de 5\% quando $LR > 3,84$.
\end{resultado}

O numerador é a verossimilhança de $x$ violações em $T$ dias se a probabilidade for a prometida, e o denominador é a mesma verossimilhança no valor de $\hat{p}$ que melhor explica os dados. A razão é próxima de 1, e $LR$ próxima de zero, quando $\hat p$ está perto de $p$. Para um ano de 250 pregões a 99\%, esperam-se 2,5 violações. Com seis, $LR = 3,56$ e o modelo ainda passa; com sete, $LR = 5,50$ e ele é rejeitado.\nota{O Comitê de Basileia usa a mesma contagem em 250 dias a 99\% e só pune o excesso: até quatro violações, zona verde; de cinco a nove, amarela; dez ou mais, vermelha. \citet{jorion2007} expõe a regra.} No outro extremo, nenhuma violação em 250 dias dá $LR = 5,03$ (as três contas estão em \arq{cap05/exemplos.py}) e a aproximação qui-quadrado rejeitaria o modelo, embora um modelo correto passe um ano inteiro sem violação com probabilidade $0,99^{250} = 8\%$. Com tão poucos eventos esperados a aproximação é grosseira, e o teste tem pouco poder para distinguir um modelo bom de um medíocre.

O Código~\ref{cod:c05-kupiec} testa três modelos nos 3624 dias de 2011 a 2025 para os quais há uma janela completa de 250 observações: o histórico com janela de 250 dias, o normal com média e desvio-padrão da mesma janela, e o normal com volatilidade EWMA\index{EWMA}, a média móvel exponencial dos retornos ao quadrado, $\sigma_t^2 = 0,94\,\sigma_{t-1}^2 + 0,06\,R_{t-1}^2$, o parente mais simples dos modelos de volatilidade condicional do Capítulo~\ref{cap:series}. É a única previsão de volatilidade usada neste capítulo; o VaR com a previsão de um GARCH fica para o último exercício.

\lstinputlisting[style=py, caption={Backtest de três modelos de VaR e teste de Kupiec (\arq{cap05/l05\_kupiec.py})}, label=cod:c05-kupiec]{scripts/cap05/l05_kupiec.py}
\begin{lstlisting}[style=saida]
95% histórico 250d 199 em 3624 (5.49%)  LR =  1.79  p = 0.1814
95% normal 250d    181 em 3624 (4.99%)  LR =  0.00  p = 0.9878
95% normal EWMA    162 em 3624 (4.47%)  LR =  2.22  p = 0.1365
99% histórico 250d  46 em 3624 (1.27%)  LR =  2.45  p = 0.1178
99% normal 250d     47 em 3624 (1.30%)  LR =  2.95  p = 0.0858
99% normal EWMA     58 em 3624 (1.60%)  LR = 11.16  p = 0.0008
\end{lstlisting}

A função usa \cod{xlogy}, que calcula $x \ln y$ e devolve zero quando $x = 0$, para que os casos de nenhuma violação e de violação todos os dias não quebrem a conta. A linha \cod{var.shift(1)} é a que impede o modelo de enxergar a perda que ele deveria prever; sem ela o backtest usa informação do futuro, o \ingles{look-ahead} definido no Capítulo~\ref{cap:dados}. A 95\% os três modelos passam. A 99\% o EWMA normal é rejeitado: ele acompanha bem o nível da volatilidade, mas o quantil normal de 2,33 desvios-padrão é pequeno demais para retornos de cauda pesada, e sobram violações.

O histórico passa no teste nos dois níveis, e a Figura~\ref{fig:c05-backtest} mostra por que isso não basta. O teste de Kupiec só olha para o total. As 46 violações do VaR histórico de 99\% vêm em blocos: oito em menos de três semanas em fevereiro e março de 2020, quando a janela de 250 dias ainda refletia um mercado calmo, e nenhuma em anos inteiros. Um modelo correto deveria errar com a mesma probabilidade todos os dias, sem que uma violação anunciasse a seguinte. O teste de \citet{christoffersen1998} acrescenta ao de Kupiec exatamente essa verificação de independência.

\figuralarga[b]{cap05/backtest}{Backtest do VaR histórico de 99\% com janela de 250 dias, de 2011 a 2025. O total de 46 violações em 3624 dias passa no teste de Kupiec, mas elas não se distribuem ao acaso: nove ocorrem em 2020, oito delas entre 26 de fevereiro e 16 de março, enquanto 2021 e 2023 não têm nenhuma.}{fig:c05-backtest}

\begin{armadilha}{Passar no teste de Kupiec não valida o modelo}
A contagem certa de violações no total é compatível com um VaR alto demais durante anos e baixo demais nas semanas em que mais se precisa dele. Olhe sempre o gráfico das violações no tempo. Blocos de violações indicam que o modelo demora a reagir a mudanças de volatilidade.
\end{armadilha}

\section{O que o VaR não vê}

O VaR tem dois defeitos que não dependem do método de cálculo. O primeiro já apareceu: ele informa onde a cauda começa e ignora o que há dentro dela. Um operador avaliado apenas pelo VaR tem incentivo para montar posições que rendem um pouco quase sempre e perdem muito com probabilidade menor que $1-\alpha$.

O segundo defeito é menos evidente e aparece quando se juntam carteiras. \citet{artzner1999} propuseram quatro propriedades que uma medida de risco $\rho$ deveria ter para merecer o nome. Na convenção de perdas deste capítulo, elas são:

\begin{itemize}
  \item \textbf{Monotonicidade.} Se $L_1 \le L_2$ em todos os cenários, então $\rho(L_1) \le \rho(L_2)$.
  \item \textbf{Invariância por translação.} Para qualquer constante $c$, $\rho(L + c) = \rho(L) + c$: somar uma perda certa $c$ aumenta o risco em $c$, e acrescentar caixa o reduz no mesmo valor.
  \item \textbf{Homogeneidade positiva.} Para $\lambda > 0$, $\rho(\lambda L) = \lambda\,\rho(L)$: dobrar a posição dobra o risco.
  \item \textbf{Subaditividade.} $\rho(L_1 + L_2) \le \rho(L_1) + \rho(L_2)$: juntar duas carteiras não pode criar risco.
\end{itemize}

Uma medida que satisfaz as quatro é chamada de \termo{medida coerente de risco}. O VaR satisfaz as três primeiras e, em geral, falha na quarta.

\begin{exemplo}{Diversificar aumenta o VaR}
Dois títulos de emissores diferentes têm, cada um, 4\% de probabilidade de calote, com perda de R\$~100 nesse caso e zero no contrário. Os calotes são independentes. Para um título sozinho, $P(L \le 0) = 96\% \ge 95\%$, logo $\VaR_{0,95} = 0$. Para a carteira com os dois, a probabilidade de pelo menos um calote é $1 - 0,96^2 = 7,84\%$, maior que 5\%, e portanto $\VaR_{0,95} = 100$. O VaR da soma (100) é maior que a soma dos VaR ($0 + 0$): a medida pune a diversificação e premia a concentração em um único emissor.
\end{exemplo}

A falha não ocorre em qualquer situação. Para carteiras lineares de ativos com distribuição conjunta elíptica, família que inclui a normal e a t multivariadas, o VaR com $\alpha \ge 50\%$ é subaditivo \citep{mcneil2015}; é o caso do exemplo de dois ativos da Seção~\ref{sec:c05-param}. O problema surge com distribuições muito assimétricas ou discretas, típicas de crédito e de opções vendidas, e basta para que limites de risco definidos por VaR em cada mesa não garantam o VaR do banco inteiro.

\section{Expected shortfall (CVaR)}

A medida que corrige os dois defeitos responde à pergunta que o VaR deixa em aberto: nos dias em que a perda supera o VaR, de quanto ela é em média?

\begin{definicao}{Expected shortfall}
O \termo{expected shortfall} (ES) de nível $\alpha$, que este livro escreve $\CVaR_\alpha$, é a média dos VaR de todos os níveis acima de $\alpha$:
\begin{equation}
  \CVaR_\alpha(L) = \frac{1}{1-\alpha}\int_\alpha^1 \VaR_u(L)\dd u .
  \label{eq:c05-es}
\end{equation}
Se $L$ tem distribuição contínua, isso equivale à perda esperada condicional a superar o VaR: $\CVaR_\alpha(L) = \E[L \mid L \ge \VaR_\alpha(L)]$.
\end{definicao}

O nome \termo{CVaR} (\ingles{conditional value at risk}) é o usado por \citet{rockafellar2000}, que mostraram como minimizar a medida por programação linear; os dois nomes designam a mesma medida, e este capítulo usa ES e CVaR como sinônimos. Por construção $\CVaR_\alpha \ge \VaR_\alpha$, e o ES é uma medida coerente \citep{acerbi2002}. No exemplo dos dois títulos, o ES de 95\% de um título é 80, a soma dos dois é 160, e o ES da carteira é 103,2 (as contas estão em \arq{cap05/exemplos.py}): a diversificação volta a reduzir o risco.

\begin{intuicao}
O VaR de 95\% é a melhor das perdas entre os 5\% piores dias. O ES de 95\% é a média de todas elas. Na Figura~\ref{fig:c05-perdas}, o VaR é a borda esquerda da região em ocre e o ES é o seu centro de gravidade, que fica em R\$~34,9 mil.
\end{intuicao}

\subsection{Fórmulas fechadas}

\begin{resultado}{ES normal}
Se $L \sim N(\mu_L, \sigma_L^2)$, então
\begin{equation}
  \CVaR_\alpha(L) = \mu_L + \sigma_L\, \frac{\varphi(z_\alpha)}{1-\alpha} ,
  \label{eq:c05-es-normal}
\end{equation}
em que $\varphi$ é a densidade da normal padrão.
\end{resultado}

A demonstração usa a propriedade $\varphi'(x) = -x\,\varphi(x)$ da densidade normal e fica como exercício. Os multiplicadores de $\sigma_L$ são 2,063 a 95\%, 2,338 a 97,5\% e 2,665 a 99\%.

Para a t de Student com $\nu > 1$ graus de liberdade, locação $\mu_L$ e escala $s$, a fórmula correspondente é \citep{mcneil2015}
\begin{equation}
  \CVaR_\alpha(L) = \mu_L + s\, \frac{f_\nu(t_{\nu,\alpha})}{1-\alpha} \cdot \frac{\nu + t_{\nu,\alpha}^2}{\nu - 1} ,
  \label{eq:c05-es-t}
\end{equation}
com $f_\nu$ a densidade da t padrão. O estimador histórico é a média das perdas da amostra que igualam ou superam o VaR histórico.

\needspace{9\baselineskip}
\lstinputlisting[style=py, caption={VaR e ES pelos três caminhos (\arq{cap05/l06\_es.py})}, label=cod:c05-es]{scripts/cap05/l06_es.py}
\needspace{11\baselineskip}
\begin{lstlisting}[style=saida]
nível   método      VaR       ES    ES/VaR
 95.0%  histórico   22,914   34,929   1.52
 95.0%  normal      25,456   32,152   1.26
 95.0%  t           23,156   34,354   1.48
 97.5%  histórico   29,595   44,412   1.50
 97.5%  normal      30,506   36,561   1.20
 97.5%  t           30,095   42,521   1.41
 99.0%  histórico   37,561   60,794   1.62
 99.0%  normal      36,377   41,808   1.15
 99.0%  t           40,258   54,853   1.36
\end{lstlisting}

A coluna ES/VaR é a que mais ensina. Sob normalidade a cauda decai tão depressa que o ES fica pouco acima do VaR, e a razão diminui com o nível de confiança: 1,26, 1,20, 1,15. Nos dados ela fica em torno de 1,5 e não diminui. A 99\%, o VaR normal e o histórico diferem em 3\%, e alguém poderia concluir que a hipótese normal é inofensiva. No ES de 99\% a diferença é de R\$~41,8 mil para R\$~60,8 mil, 45\% a mais. A normal erra pouco a posição do quantil e erra muito o que vem depois dele; a t chega mais perto, e ainda assim fica abaixo do histórico.

\subsection{Do VaR de 99\% ao ES de 97,5\%}

Na revisão das regras de capital para risco de mercado conhecida como FRTB\index{FRTB} (\ingles{Fundamental Review of the Trading Book}), o Comitê de Basileia substituiu o VaR de 99\% pelo ES de 97,5\% nos modelos internos dos bancos \citep{bcbs2019}. A escolha do nível tem uma lógica que a Equação~\eqref{eq:c05-es-normal} revela: sob normalidade, o ES de 97,5\% vale $2,338\,\sigma_L$ e o VaR de 99\% vale $2,326\,\sigma_L$, praticamente o mesmo número. Para uma carteira de cauda normal a troca não altera a exigência; para caudas pesadas, ela a aumenta. Na carteira do capítulo, o ES histórico de 97,5\% é de R\$~44,4 mil, contra um VaR histórico de 99\% de R\$~37,6 mil.

O ES tem um custo estatístico: é uma média de observações raras, e por isso sua estimativa é mais instável que a do VaR de mesmo nível (Seção~\ref{sec:c05-cauda}) e mais difícil de conferir em um backtest.

\section{Drawdown}
\label{sec:c05-drawdown}

VaR e ES olham para um dia de cada vez. O investidor sente o risco de outra forma: pela queda acumulada desde o melhor momento da carteira e pelo tempo que ela leva para voltar até lá. Uma sequência de trinta perdas pequenas, nenhuma delas uma violação de VaR, pode custar mais do que um dia de pânico.

\begin{definicao}{Drawdown}
Seja $V_t$ o valor da carteira (ou da cota de um fundo) no instante $t$. O máximo atingido até $t$, $M_t = \max_{s \le t} V_s$, é o \termo{high-water mark}. O \termo{drawdown} em $t$ é a distância percentual até ele,
\begin{equation}
  D_t = \frac{V_t}{M_t} - 1 \le 0 ,
  \label{eq:c05-dd}
\end{equation}
e o \termo{drawdown máximo} de um período é o menor valor de $D_t$ nesse período.
\end{definicao}

Três números descrevem um episódio de drawdown: a profundidade (o valor de $D_t$ no vale), a duração da queda, do pico ao vale, e o \termo{tempo de recuperação}, do vale até a cota voltar ao pico anterior. A soma das duas durações, de um pico ao seguinte, é o \termo{tempo submerso}. O Código~\ref{cod:c05-dd} calcula os quatro para a \cod{cota} do Código~\ref{cod:c05-carteira}, com os dezesseis anos de dados, de 2010 a 2025, e supondo que os pesos iguais são refeitos todos os dias.

\lstinputlisting[style=py, caption={Série de drawdown, drawdown máximo e recuperação (\arq{cap05/l07\_drawdown.py})}, label=cod:c05-dd]{scripts/cap05/l07_drawdown.py}
\begin{lstlisting}[style=saida]
drawdown máximo: -42.8%
pico 05/05/2015  vale 26/01/2016  recuperação 15/08/2016
queda: 179 pregões   recuperação: 139   tempo submerso: 318
retorno a.a.: 16.1%   Calmar: 0.38
fração dos dias abaixo do pico: 92.1%
\end{lstlisting}

O método \cod{cummax()} devolve o máximo acumulado e resolve a conta em uma linha. O pior episódio foi a queda de 42,8\% entre maio de 2015 e janeiro de 2016, durante a recessão brasileira: 179 pregões de descida e mais 139 para recuperar, 318 de tempo submerso (o código conta as datas da cota, que no episódio coincidem com os pregões). O retorno é anualizado pela convenção de 252 dias (Seção~\ref{sec:anualizar}). A última linha merece atenção. A carteira multiplicou o capital por mais de dez no período e, mesmo assim, passou 92\% dos dias abaixo de algum pico anterior.

\figuralarga[t]{cap05/underwater}{Cota da carteira, \ingles{high-water mark} e drawdown, de 2010 a 2025. Os dois maiores episódios têm profundidade parecida e formatos opostos: o de 2015--2016 levou 179 pregões do pico ao vale; o de 2020 levou 18 pregões para cair 39,9\% e 86 para recuperar.}{fig:c05-underwater}

O gráfico de baixo da Figura~\ref{fig:c05-underwater} é chamado de \ingles{underwater}\index{underwater}, porque mostra quanto tempo a carteira passa ``debaixo d'água''. Ele deixa ver o que o número isolado esconde: 2015--2016 e 2020 têm quase a mesma profundidade, mas o primeiro foi uma erosão de nove meses e o segundo um choque de menos de um mês.

A recuperação exige mais do que a queda tirou, pela aritmética da armadilha sobre somar retornos do Capítulo~\ref{cap:retornos}: depois de um drawdown $D$, o retorno necessário para voltar ao pico é $1/(1+D) - 1$. Para $D = -42,8\%$, são $+74,8\%$.

O \termo{índice de Calmar} resume a relação entre o que a carteira rendeu e o pior momento que impôs: divide o retorno anualizado pelo módulo do drawdown máximo do mesmo período. A medida foi proposta por \citet{young1991}, que usava o retorno anual médio e o drawdown máximo dos últimos 36 meses. Para a carteira, $16,1\% / 42,8\% = 0,38$: cada ponto de retorno anual veio acompanhado de cerca de 2,7 pontos de drawdown máximo. O Capítulo~\ref{cap:backtest} retoma o Calmar ao lado das demais métricas de desempenho.

\begin{armadilha}{O drawdown máximo é uma única observação}
O drawdown máximo depende de um só episódio e do tamanho da amostra: quanto mais longo o histórico, maior tende a ser o pior episódio encontrado, e por isso não se comparam drawdowns de fundos com idades diferentes. Há ainda um viés nesta carteira em particular: as cinco ações foram escolhidas hoje, entre empresas que sobreviveram e prosperaram.
\end{armadilha}

\section{Métricas de cauda e o tamanho do erro}
\label{sec:c05-cauda}

Além de VaR, ES e drawdown, relatórios de risco costumam trazer um conjunto de estatísticas descritivas da cauda. O Código~\ref{cod:c05-cauda} calcula as mais comuns.

\needspace{14\baselineskip}
\lstinputlisting[style=py, caption={Métricas descritivas de cauda (\arq{cap05/l08\_cauda.py})}, label=cod:c05-cauda]{scripts/cap05/l08_cauda.py}
\begin{lstlisting}[style=saida]
assimetria: -0.53
excesso de curtose: 10.43
razão de cauda q95/|q05|: 1.07
pior dia: -15.0% em 12/03/2020
pior mês: -22.6% em 03/2020
pior dia em desvios-padrão: -9.4
probabilidade normal de um dia assim: 1.9e-21
\end{lstlisting}

A assimetria e a curtose foram definidas no Capítulo~\ref{cap:estatistica}. Assimetria negativa indica que os desvios grandes ocorrem mais para o lado das perdas, e um excesso de curtose de 10 indica caudas muito mais pesadas que as da normal, cujo excesso é zero. A \termo{razão de cauda} divide o quantil de 95\% dos retornos pelo módulo do quantil de 5\%: acima de 1, os ganhos típicos dos dias bons superam as perdas típicas dos dias ruins. Aqui ela vale 1,07 enquanto a assimetria é negativa, e não há contradição: a razão olha para os quantis de 5\% e a assimetria é dominada pelos poucos dias extremos. Pior dia e pior mês são os números mais fáceis de comunicar e os mais frágeis, porque cada um é uma observação.

As duas últimas linhas dão a medida do fracasso da normal. O pior dia ficou a 9,4 desvios-padrão da média, evento ao qual a normal atribui probabilidade da ordem de $10^{-21}$.

\begin{repositorio}
A pasta \arq{aula5/} traz os três tipos de VaR para VALE3 (\arq{Var\_parametrico.py}, \arq{Var\_historico.py}, \arq{Var\_MonteCarlo.py}), notebooks de CVaR e de drawdown máximo, e o cálculo de assimetria e curtose com laços (\arq{sk.py}, \arq{kurtose.py}). Os scripts de VaR informam o resultado como log-retorno negativo; refaça-os com perdas positivas em reais e confira com os códigos deste capítulo.
\end{repositorio}

\subsection{Por que toda medida de cauda é imprecisa}

Uma medida de cauda é estimada com as poucas observações que caem na cauda. Em um ano de 250 pregões, o VaR de 95\% se apoia em 12 ou 13 perdas, o de 99\% em duas ou três, e o ES de 99\% é a média dessas duas ou três. A Figura~\ref{fig:c05-incerteza} mede a consequência por \ingles{bootstrap}\index{bootstrap}: sorteiam-se, com reposição, amostras de 250 dias das perdas da carteira, recalcula-se cada medida e observa-se quanto ela varia de uma amostra para outra.

\figura[t]{cap05/incerteza}{Intervalos de 90\% de quatro medidas de risco em 5000 reamostragens das perdas diárias da carteira (2015--2024). Com 250 dias, o VaR de 99\% pode sair entre R\$~29 mil e R\$~54 mil e o ES de 99\% entre R\$~34 mil e R\$~93 mil. Quanto mais para dentro da cauda, maior o erro.}{fig:c05-incerteza}

Com um ano de dados, o desvio-padrão é estimado com um intervalo de largura igual a um terço do seu valor; o VaR de 99\%, com um intervalo de 70\% do valor; o ES de 99\%, com um intervalo maior que a própria estimativa. Com dez anos os intervalos encolhem, mas o do ES de 99\% ainda vai de R\$~48,5 mil a R\$~73,5 mil. E a figura é otimista: o sorteio trata os dias como independentes, quando na realidade os dias ruins vêm agrupados e a informação contida na amostra é menor do que o número de observações sugere.

A lição prática é desconfiar de casas decimais. Um VaR de 99\% informado como R\$~37.561 tem bem menos precisão do que aparenta: mesmo com dez anos de dados, o intervalo de 90\% do bootstrap vai de R\$~35 mil a R\$~39 mil, e o segundo algarismo já é incerto. A diferença entre os métodos da Tabela~\ref{tab:c05-comparacao} é menor que o erro de estimação de qualquer um deles com um ano de dados, e por isso a escolha da janela e a reação do modelo a mudanças de volatilidade pesam mais do que a escolha entre fórmulas. \citet{mcneil2015} tratam da teoria dos valores extremos, o conjunto de técnicas estatísticas desenvolvido para extrapolar a cauda além do que a amostra mostra.

\begin{resumo}
  \item O VaR de nível $\alpha$ é o quantil $\alpha$ das perdas em um horizonte dado. Ele marca onde a cauda começa e não é a perda máxima.
  \item O VaR paramétrico normal é $\mu_L + \sigma_L z_\alpha$, com $\sigma_L = V\sqrt{\vetor{w}\T\symbf{\Sigma}\vetor{w}}$. A normal tende a superestimar o VaR de 95\% e a subestimar o de 99\%; com a t de Student, a escala precisa do fator $\sqrt{(\nu-2)/\nu}$.
  \item O VaR histórico dispensa hipóteses de distribuição e depende da janela. O Monte Carlo só acrescenta algo quando a carteira é não linear ou o modelo é dinâmico.
  \item Um modelo de VaR se testa contando violações fora da amostra: o teste de Kupiec confere o total, e o gráfico no tempo mostra se vêm em blocos.
  \item O VaR não é subaditivo em geral. O \ingles{expected shortfall}, a perda média além do VaR, é coerente e sensível à forma da cauda, e foi adotado pela regulação bancária a 97,5\%.
  \item O drawdown mede a queda acumulada desde o pico. Toda medida de cauda, do VaR de 99\% ao drawdown máximo, se apoia em poucas observações e tem erro de estimação grande.
\end{resumo}

\begin{exercicios}
  \item \nivel{básico} Uma carteira de R\$~500 mil tem volatilidade diária de 1,2\%. Supondo normalidade e média zero, calcule à mão o VaR de um dia a 95\% e a 99\%, o VaR de dez dias a 99\% pela raiz do tempo e o ES de um dia a 97,5\%.
  \item \nivel{básico} Um relatório informa: ``VaR de um dia a 99\%: R\$~40 mil''. Verdadeiro ou falso, e por quê: (a) a carteira nunca perde mais de R\$~40 mil em um dia; (b) em dois ou três dias por ano a perda deve superar R\$~40 mil; (c) nos dias de violação a perda média é de R\$~40 mil.
  \item \nivel{intermediário} Um modelo de VaR de 95\% teve 36 violações em 500 dias. Calcule a estatística de Kupiec pela Equação~\eqref{eq:c05-kupiec} e decida ao nível de 5\%. Para que números de violações, acima e abaixo do esperado (25), o modelo é rejeitado?
  \item \nivel{intermediário} Demonstre a Equação~\eqref{eq:c05-es-normal} e encontre numericamente o nível $\alpha$ em que o ES normal iguala o VaR normal de 99\%.
  \item \nivel{intermediário} Adapte \arq{cap05/carteira.py} para uma posição de R\$~100 mil em \cod{PETR4.SA}. Calcule o VaR e o ES históricos de 99\% de 2015 a 2024 e liste os três drawdowns mais profundos de 2010 a 2025, com profundidade, duração da queda e tempo de recuperação.
  \item \nivel{desafio} A expansão de Cornish--Fisher corrige o quantil normal pela assimetria $S$ e pelo excesso de curtose $K - 3$ da perda (notação do Capítulo~\ref{cap:estatistica}): $\tilde{z}_\alpha = z_\alpha + (z_\alpha^2 - 1)S/6 + (z_\alpha^3 - 3z_\alpha)(K-3)/24 - (2z_\alpha^3 - 5z_\alpha)S^2/36$, com $\VaR_\alpha = \mu_L + \sigma_L \tilde{z}_\alpha$. Aplique-a às perdas diárias e mensais da carteira e compare com o VaR histórico. O que acontece quando o excesso de curtose é da ordem de 10?
  \item \nivel{desafio} Acrescente ao Código~\ref{cod:c05-kupiec} um modelo com volatilidade prevista por um GARCH(1,1) com inovações t (Capítulo~\ref{cap:series}) e quantil da t com a escala correta. Ele passa no teste de Kupiec a 99\%? As violações de 2020 continuam agrupadas?
\end{exercicios}

\begin{leituras}
  \item \citet{jorion2007}, a referência clássica sobre VaR, para os métodos de cálculo, o backtest e a simulação de Monte Carlo.
  \item \citet{hull2018}, para carteiras com opções e para a regulação de Basileia.
  \item \citet{mcneil2015}, para o tratamento rigoroso de VaR, ES, medidas coerentes e teoria dos valores extremos.
  \item \citet{artzner1999}, o artigo que formulou os axiomas de coerência, e \citet{acerbi2002}, que demonstram que o ES os satisfaz.
\end{leituras}
